\section{Methodology and Proposed System}
\label{sec:method}

\subsection{System Overview and Problem Formulation}
\label{sec:model}

\Cref{fig:arch} shows the architecture: a checkpointed LangGraph state machine drives a Scan$\to$Reason$\to$Act$\to$Observe loop behind three typed ports. The remainder of this subsection formalises what the loop explores and what it checks.

\begin{figure*}[t]
\centering
\resizebox{\textwidth}{!}{%
\begin{tikzpicture}[
  font=\footnotesize,
  box/.style={draw=acc,thick,rounded corners=2pt,fill=acc!6,minimum height=9mm,
              minimum width=24mm,align=center},
  port/.style={draw=mute,dashed,rounded corners=2pt,inner sep=5pt},
  data/.style={draw=mute,fill=white,rounded corners=1pt,align=center,font=\scriptsize},
  arr/.style={-{Stealth[length=2mm]},thick,draw=black!70},
  lbl/.style={font=\scriptsize,fill=white,inner sep=1pt}]

  \node[box] (scan)    {\textbf{scan}\\fingerprint $\phi$ · persist $v$\\enumerate actions};
  \node[box,right=17mm of scan] (reason) {\textbf{reason}\\$\hat\Sigma(v)$ · audit\\$\mathit{Viol}_\mathcal F\cup\mathit{Viol}_\mathcal R$};
  \node[box,right=17mm of reason] (act) {\textbf{act}\\pop frontier · replay path\\claim $(v,a)$ · execute};
  \node[box,right=12mm of act] (obs) {\textbf{observe}\\persist edge $(v,a,v')$\\back-edge kept};

  \draw[arr] (scan) -- (reason);
  \draw[arr] (reason) -- node[lbl,above=1mm,align=center]{clean /\\violated} (act);
  \draw[arr] (act) -- (obs);
  \draw[arr] (obs.south) -- ++(0,-5mm) -| node[lbl,pos=0.25,above]{progressed} (scan.south);
  \draw[arr] (obs.north) -- ++(0,4mm) -| node[lbl,pos=0.25,above]{skipped (unreplayable / stale)} (act.north);
  \node[right=4mm of obs,font=\scriptsize,align=left] (end) {\textsc{end}\\(frontier exhausted,\\ state cap, stop)};
  \draw[arr,dashed] (obs) -- (end);

  \node[data,below=17mm of scan]   (crawl) {\textbf{CrawlerPort}\\Playwright: \texttt{open}, \texttt{act},\\\texttt{enumerate\_actions}};
  \node[data,below=17mm of reason] (perc)  {\textbf{PerceptionPort}\\DOM/AX correlator · VLM\\\texttt{analyze}, \texttt{audit}};
  \node[data,below=17mm of act]    (graph) {\textbf{GraphPort}\\claim set $C$ (Redis)\\cyclic graph (Neo4j)};
  \node[data,below=17mm of obs]    (pol)   {\textbf{Policy}\\English $\to$ clauses\\$(\sigma,\rho,\pi)$};

  \begin{scope}[on background layer]
    \node[port,fit=(crawl)(perc)(graph)(pol),label={[font=\scriptsize,text=mute]below:typed, frozen contracts (\texttt{contracts.py}) --- injected, never imported}] {};
  \end{scope}
  \foreach \a/\b in {scan/crawl,reason/perc,act/graph} \draw[-,thin,draw=mute,dotted] (\a) -- (\b);
  \draw[-,thin,draw=mute,dotted] (obs) -- (pol);
\end{tikzpicture}}
\caption{\sys{} architecture. A checkpointed LangGraph state machine of four
nodes and two conditional edges drives the Scan$\to$Reason$\to$Act$\to$Observe
loop. All I/O goes through three typed ports, so every component can be
replaced by a deterministic fake for testing and by the live implementation in
production without changing the loop.}
\label{fig:arch}
\end{figure*}


\subsubsection{Applications, roles, and observations}

We model a web application under test (AUT), as seen by a fixed role $r$, as a
labelled transition system $M_r=(Q,A,\delta,q_0)$ over concrete UI
configurations $Q$, actions $A$, and a partial transition function
$\delta:Q\times A\rightharpoonup Q$. The auditor never sees $Q$ directly. A
\emph{capture} $\kappa(q)=(u,d,t,s)$ bundles the URL, serialized DOM,
accessibility (AX) tree, and an optional screenshot. \sys{} pins one role per
run; multi-role coverage is several runs, and comparing roles is a
reporting-layer concern (\cref{sec:discussion}).

\begin{definition}[State abstraction]
A \emph{fingerprint} is $\phi = H\circ N$, where $N$ normalizes a capture by
removing volatile content (session and CSRF tokens, nonces, long numeric
timestamps, generated hex identifiers, per-session CDP node ids, transient
focus), and $H$ is SHA-256. Two configurations are \emph{the same abstract
state} iff $\phi(\kappa(q))=\phi(\kappa(q'))$.
\end{definition}

Normalization has two failure modes. If it is too sensitive, the model
explodes into one node per page load. If it is too coarse, a security-relevant
change (for example, a destructive button that appears) folds into an existing
state and is never audited. \Cref{sec:eval:fp} measures both.

\subsubsection{The exploration graph is cyclic}

\begin{definition}[Exploration graph]
The \emph{exploration graph} is a directed multigraph $G=(V,E)$ with
$V\subseteq\operatorname{img}\phi$ and $E\subseteq V\times A\times V$, where
$(v,a,v')\in E$ records that executing $a$ in $v$ was observed to lead to $v'$.
An edge is a \emph{back-edge} if $v'$ was already discovered when it was added.
\end{definition}

Model-inference crawlers often treat revisits as redundant and drop them,
leaving a tree or DAG. We treat back-edges as \emph{evidence}. ``Log out
returns to login'' and ``Admin returns to Dashboard'' are facts about the
AUT's navigation structure, and a path-sensitive audit can only reason over
them if they are kept. \sys{} therefore never prunes an edge. Loop prevention
lives entirely in a \emph{claim set} $C\subseteq V\times A$.

\begin{definition}[Claim-before-execute]
Action $a$ is executed from $v$ only if $(v,a)\notin C$, and $(v,a)$ is inserted
into $C$ \emph{before} $a$ is sent to the browser. The edge $(v,a,v')$ is
persisted only \emph{after} $a$ succeeds.
\end{definition}

\begin{proposition}[Termination and at-most-once per pair]
\label{prop:term}
Suppose $|V|\le n$, which is enforced by a state cap, and each state enables
at most $k$ actions. Then exploration performs at most $nk$ \emph{frontier}
executions and terminates. No claimed pair $(v,a)$ is executed as a frontier
action twice, including across crash and resume.
\end{proposition}
\begin{proof}[Proof sketch]
Every frontier execution inserts a distinct pair into $C$ first, and
$|C|\le nk$. On resume, a pair that is in $C$ but has no edge is treated as
done and never re-fired. Execution strictly follows insertion, so a crash
cannot cause a second execution of the same pair.
\end{proof}

Most AUTs are not idempotent: re-submitting a form changes the application
being audited. Reaching a queued action with a single browser, however, means
\emph{replaying} its path from the seed, and replay re-fires the actions on
that path. \Cref{prop:term} is therefore a guarantee about \emph{pairs}, not
about \emph{effects}. \sys{} narrows the gap by kind: only
$\textsf{ReplaySafe}=\{\texttt{navigate},\texttt{click},\texttt{back}\}$ may be
re-fired. A candidate whose path passes through a \texttt{fill},
\texttt{select}, or \texttt{submit} is recorded as skipped, with a reason, and
its subtree is reported as unexplored instead of being reached by quietly
re-submitting a form. The residual assumption is that \emph{clicks are
navigational}. A click that mutates state on a replayed path (for example, a
one-click ``Delete'') \emph{is} re-fired. We treat this as an open limitation
(\cref{sec:discussion}), not a solved problem. Action identity is
\emph{content-addressed}:
$\mathit{id}(a)=H(\mathit{role}\,\|\,\mathit{name}\,\|\,\mathit{selector})$, not
a DOM index. Replay therefore survives DOM reordering, and the same ``Log out''
control keeps one identity everywhere, while deduplication stays scoped to
$(v,a)$.

\subsubsection{Policies as two set operations}

A policy is written in English by a QA engineer and compiled into
machine-checkable clauses $c=(\sigma_c,\rho_c,\pi_c)$: a \emph{subject}
$\sigma_c$ (a UI element kind or capability), a role scope $\rho_c$ (empty
means all roles), and a polarity $\pi_c\in\{\textsf{must\_exist},
\textsf{must\_not\_exist}\}$. This splits the policy into forbidden clauses
$\Forb$ and required clauses $\Req$. Perception maps each state to a set of
grounded subjects $\Sig(v)$ (element tags, accessible names, and capabilities).
For a role $r$, let $\Forb_r$ and $\Req_r$ be the clauses whose scope admits
$r$.

\begin{definition}[Negation-aware audit]
\begin{align}
\mathit{Viol}_\Forb(v) &= \{\,c\in\Forb_r : \sigma_c\in\Sig(v)\,\}
  && \text{(forbidden present)}\label{eq:forb}\\
\mathit{Viol}_\Req(v)  &= \{\,c\in\Req_r : \sigma_c\notin\Sig(v)\,\}
  && \text{(required absent)}\label{eq:req}
\end{align}
The audit of $v$ is $\mathit{Viol}_\Forb(v)\cup\mathit{Viol}_\Req(v)$, and each
violation is tagged with the operation that produced it.
\end{definition}

\Cref{eq:forb} is an intersection ($S\cap\Forb$) and \cref{eq:req} is a
difference ($\Req\setminus S$). No single intersection can express both, so an
intersection-only oracle is structurally unable to report a missing element
(\cref{sec:eval:rq3}).

\subsubsection{Error-direction asymmetry under imperfect perception}
\label{sec:model:asym}

In a real system, $\Sig(v)$ is replaced by a perception estimate $\hat\Sig(v)$
from a DOM/AX correlator or a vision-language model. The two clause types react
to perception errors in \emph{opposite} directions.

\begin{proposition}[Error-direction asymmetry]
\label{prop:asym}
Let $\widehat{\mathit{Viol}}$ denote the audit computed on $\hat\Sig$.
\begin{enumerate}
  \item If perception \emph{under-reports} ($\hat\Sig(v)\subseteq\Sig(v)$), then
    $\widehat{\mathit{Viol}}_\Forb(v)\subseteq\mathit{Viol}_\Forb(v)$ and
    $\widehat{\mathit{Viol}}_\Req(v)\supseteq\mathit{Viol}_\Req(v)$. Forbidden
    clauses can only be missed, and required clauses can only raise false
    alarms.
  \item If perception \emph{over-reports} ($\hat\Sig(v)\supseteq\Sig(v)$, for
    example by hallucinating a subject), both inclusions reverse.
\end{enumerate}
\end{proposition}
\begin{proof}
\Cref{eq:forb} is monotone in $\Sig$ and \cref{eq:req} is antitone in $\Sig$.
\end{proof}

\begin{corollary}
No uniform perception bias is safe for both clause types. A detector tuned to
avoid false alarms on forbidden clauses (conservative, under-reporting) creates
false alarms on required clauses, and the reverse also holds.
\end{corollary}

This matters for multimodal perception. Vision-language models answer negated
queries close to chance~\cite{alhamoud2025negation}, and their compositional
understanding is known to be shallow~\cite{yuksekgonul2023aro,ma2023crepe}.
Those studies use natural images, not UI screenshots, so we cite them as
motivation rather than as a measurement of UI perception. Whatever a
perceiver's error rate, \cref{prop:asym} shows that a single yes/no ``is
$\sigma$ on the screen?'' query has a polarity-dependent risk profile. This motivates the design
choices in \cref{sec:system:negation}: evidence-carrying verdicts, a
deterministic DOM/AX channel that cross-checks the VLM, and (proposed) a
three-valued verdict that abstains instead of guessing when the two channels
disagree. \Cref{sec:eval:rq2} observes exactly the asymmetry predicted by
\cref{prop:asym}.

\subsection{Materials, Dataset and Tools}
\label{sec:materials}

\paragraph{Tools.}
The system is written in Python (about 5.8k lines across five modules) with a Playwright crawler, a LangGraph orchestrator, Redis for the claim set, Neo4j for the cyclic graph, a FastAPI reporting service, and a VS~Code extension front-end. The vision-language channel uses Gemini~2.5 Flash behind a rate-limited provider interface; it is implemented but not evaluated here.

\paragraph{Dataset: subject application.}
\texttt{broken-admin} is a static six-page application (index, login,
register, dashboard, reports, admin) with four cycles and four planted
violations, browsed as \texttt{guest} (\cref{fig:app}). The policy has three
forbidden clauses (admin link, delete-user, debug panel) and one required
clause (skip-to-content link). V-01 is an Admin link rendered for every role;
the backend rejects the request, but the link itself is the exposure. V-02 is
a debug panel left in the production build. V-03 is a Delete-user control two
clicks from the seed. V-04 is a missing skip-to-content link, which is an
\emph{absence}.

\paragraph{Ground-truth format.}
Each test application ships three machine-readable files next to its HTML:
\texttt{app.json} (pages and controls), \texttt{policy.json} (the English
policy and its structured form), and \texttt{violations.json} (every planted
defect with its page and clause). Evaluation asserts \emph{exact} set equality
between found and planted violations. A miss is a regression, and an
undocumented finding is a false positive. A consistency test fails if the HTML
and \texttt{app.json} drift apart.

\paragraph{Reproducibility.}
All numbers are produced by
\nolinkurl{docs/demo_crawl.py} (RQ1) and
\nolinkurl{docs/report/experiments/run_experiments.py} (RQ2--RQ4). The raw
output is in \nolinkurl{docs/report/data/experiments.json}. Every run is
deterministic.

\begin{figure}[t]
\centering
\begin{tikzpicture}[font=\scriptsize,
  st/.style={draw=acc,thick,rounded corners=2pt,fill=acc!6,minimum width=15mm,minimum height=6mm},
  vio/.style={st,draw=bad,fill=bad!8},
  fe/.style={-{Stealth[length=1.8mm]},thick},
  be/.style={-{Stealth[length=1.8mm]},thick,dashed,draw=bad!80}]
  \node[st]  (idx) at (0,0)      {index};
  \node[st]  (log) at (2.7,0)    {login};
  \node[st]  (reg) at (5.6,0)    {register};
  \node[vio] (das) at (2.7,-1.7) {dashboard$^{\text{V1}}$};
  \node[vio] (rep) at (0.4,-3.3) {reports$^{\text{V2}}$};
  \node[vio] (adm) at (5.0,-3.3) {admin$^{\text{V3,V4}}$};
  \draw[fe] (idx) -- (log);
  \draw[fe] (log) to[bend left=18] (reg);
  \draw[be] (reg) to[bend left=18] (log);
  \draw[fe] ([xshift=-2mm]log.south) -- ([xshift=-2mm]das.north);
  \draw[be] ([xshift=2mm]das.north) -- node[right,font=\tiny]{log out} ([xshift=2mm]log.south);
  \draw[fe] (das) to[bend right=14] (rep);
  \draw[be] (rep) to[bend right=14] (das);
  \draw[fe] (das) to[bend left=14] (adm);
  \draw[be] (adm) to[bend left=14] (das);
  \draw[be] (rep.north) to[out=100,in=200] (log.west);
  \draw[be] (adm.north east) to[out=60,in=-30] (log.south east);
\end{tikzpicture}
\caption{Exploration graph recovered for \texttt{broken-admin} (role
\texttt{guest}): 6 states, 11 edges, of which 6 are back-edges (dashed). A
tree- or DAG-based crawler would drop every dashed edge. Superscripts mark
planted violations.}
\label{fig:app}
\end{figure}

\subsection{Algorithm and Implementation Details}
\label{sec:system}




\Cref{fig:arch} shows the architecture. The implementation consists of about
5.8k lines of Python across five modules owned by four sub-teams, plus a
FastAPI reporting service and a VS~Code extension front-end.

\subsubsection{Orchestration}
The explorer performs breadth-first search over $(v,a)$ pairs. A plain-Python
reference loop, kept as a \emph{parity oracle}, is ported to a LangGraph
\texttt{StateGraph} with four nodes (\texttt{scan}, \texttt{reason},
\texttt{act}, \texttt{observe}) and two conditional edges. The whole
exploration state (frontier, claimed pairs, discovery order, skips) is a
single checkpointed \texttt{TypedDict}, so a crashed run can resume instead of
restarting. Limits such as depth, state cap, and perception rate are
configuration rather than state, so a resumed run cannot silently run under
different rules. A separate conditional edge routes a skipped action back to
\texttt{act} and not through \texttt{scan}. Otherwise an unchanged state would
be re-audited and a phantom entry added to the discovery order, and the
discovery order is the evidence for cycles. The same test suite runs against
both implementations.

\subsubsection{Capture and actions}
A Playwright crawler captures DOM, the CDP full AX tree, and a screenshot. It
derives candidate actions from actionable AX roles (link, button, tab,
menuitem, \dots) and gives each one a content-addressed identifier. Navigation
failures and action failures raise distinct exception types, so the
orchestrator can retry the former with backoff and mark-and-skip the latter
without type-sniffing.

\subsubsection{State abstraction and persistence}
$\phi$ removes known volatile attributes (\texttt{nonce},
\texttt{data-session}, \texttt{data-csrf}, \dots), hex-like generated ids,
prefixed tokens (\texttt{sess\_}, \texttt{csrf\_}, \texttt{tok\_}), 10+-digit
timestamps, whitespace, and volatile CDP keys (\texttt{nodeId},
\texttt{backendDOMNodeId}, \texttt{parentId}, \texttt{childIds}, focus). It
then hashes URL, DOM, and canonical AX JSON together. The claim set lives in
Redis. States, edges (including back-edges), and violations are stored in
Neo4j, where the cyclic structure can be queried directly (for example, ``all
paths from the seed to a violating state'').

\subsubsection{Perception: grounding subjects in the UI}
\label{sec:system:perception}
Perception turns a capture into a \emph{Semantic UI Map}: elements with role,
accessible name, tags, selector, visibility and enabled state, plus the
capabilities the state exposes (\texttt{delete-user}, \texttt{export-data},
\dots). The baseline channel is deterministic and fully reproducible. It
correlates DOM elements and AX nodes by (i)~browser-provided identifiers,
(ii)~role plus accessible name, and (iii)~name alone when unique. Ambiguous
matches are never guessed. DOM-only elements are kept on purpose, because an
element hidden from the AX tree can still be an exposure. A vision-language
model channel (Gemini~2.5 Flash, behind a rate-limited provider interface)
reads the screenshot and AX tree. It supplies judgements the DOM cannot make,
such as a styled \texttt{<div>} with no ARIA role that \emph{looks like} a
delete button. By design it may only \emph{annotate} elements the
deterministic channel already found, adding tags and capabilities. It may never
add or remove an element, so existence and absence (the polarity that VLMs are
known to handle poorly~\cite{alhamoud2025negation}) are decided by evidence
that cannot be hallucinated. In the current codebase this channel is
implemented but not yet wired into the orchestrator, and it is not evaluated
in this paper. A run without screenshots is labelled a smoke test, not an
audit.

\subsubsection{Negation engine}
\label{sec:system:negation}
The audit implements \cref{eq:forb,eq:req} deterministically. Clause
applicability is checked against the run's role. Subjects are matched first
against capabilities, which are stronger evidence, and then against element
tags and accessible names. There is no substring matching, because an
over-eager matcher is exactly the over-reporting bias of \cref{prop:asym}.
Every violation carries its clause type, a severity, a plain-language
rationale, and \emph{evidence} (selector, text, screenshot) that a human can
check. The policy compiler that turns English sentences into clauses is
\emph{in progress}. By design it reuses the perception port's text channel, so
there is one model client and one rate limit, and it includes a confirmation
gate on which the QA engineer approves the parse before the crawl. In this
draft, policies are supplied in structured form next to their English source
(\texttt{policy.json}).

